\section{Closed-loop results: safety, recovery and compute}\label{sec:closedloop}

Offline evaluation scores one proposal per episode. A deployed screen acts inside the propose--validate--retry
loop. There its rejections create new proposals, its decisions change which feedback the agent sees, and the
cost of a skipped call depends on what replaces it. This section answers RQ4 on 400 fresh CSTR episodes per
model (100 per fault family), disjoint from all data used to build the screens.

\subsection{Policies and protocol}
Every policy was run on the same 400 episodes with the frozen probes of each model. The policies are listed
below.
\begin{itemize}
\item \emph{Always validate} (AV), the conventional loop and the reference.
\item \emph{Never validate}, which executes the first proposal.
\item \emph{Random routing}, which accepts, rejects or validates at random with the combined probe's development
rates (a control for the skip rate itself).
\item \emph{Observables}, the observable screen.
\item \emph{Combined}, the screen on observables, all internals and grounding.
\item \emph{Internals only}, the screen on all internals and grounding, without plant readings or proposed
setpoints.
\end{itemize}
The probe-based policies were also run under the first-proposal-only rule (R0) wherever R0 could change the
outcome (below). An unchecked rejection returns the fixed message ``the previous proposal was rejected by a risk
screen before validation; propose a different set of setpoints'', which contains no measured metric. A
validated failure returns the validator's measured outcome. A proposal that repeats an earlier one keeps its
earlier verdict and costs no new call. Executing a proposal means applying its setpoints to the plant from the
snapshot, with the plant's own future disturbances. The executed outcome is the simulator's verdict on that
trajectory, so the validator doubles as the plant. Its decision-time calls are charged as compute, while the
executed trajectory is not.

\subsection{Safety: trained screens versus skipping at random}
Figure~\ref{fig:plane} places every policy in the plane of validator calls avoided against failing actions
executed without validation. Table~\ref{tab:closedloop} gives the numbers. Validating every proposal recovered
34.0--53.8\% of episodes and, by construction, never executed a failing action. Never validating executed a
failing action in 63.5--79.5\% of episodes. Random routing at the screens' rates executed 107--186 failing actions
per 400 episodes. The exception was Llama-3.2-3B, whose development accept rate, and hence the random policy's,
was zero. The trained screens avoided 29--90\% of validator calls. On first proposals they behaved as they had
offline: the observable screens executed 0--13 failing actions unchecked per 400 episodes. The difference to
random routing at the same skip rates shows that the screens' safety comes from \emph{which} proposals they skip,
not from how many.

\begin{figure}[!htbp]
\centering
% Closed loop: validator calls avoided against failing actions executed without validation (400 episodes per model).
% Arrows: effect of the first-proposal-only rule (R0). Data: data/plane_*.dat (scripts/50).
\begin{tikzpicture}
\pgfplotsset{
  pAV/.style={only marks, mark=square, mark size=2.2pt, black, thick},
  pNV/.style={only marks, mark=x, mark size=3pt, black, thick},
  pRand/.style={only marks, mark=diamond*, mark size=2.6pt, cRand},
  pObs/.style={only marks, mark=*, mark size=2.1pt, cObs},
  pComb/.style={only marks, mark=triangle*, mark size=2.7pt, cComb},
  pInt/.style={only marks, mark=square*, mark size=2.0pt, cInt},
  pR0/.style={only marks, mark=o, mark size=3.6pt, cRzero, thick},
}
\begin{groupplot}[group style={group size=2 by 1, horizontal sep=14mm},
  paperaxis, width=0.47\textwidth, height=0.45\textwidth, xmin=-3, xmax=103,
  xlabel={validator calls avoided (\%)}, xtick={0,20,40,60,80,100}, xmajorgrids]
\nextgroupplot[title={\footnotesize (a) all policies}, ymin=-3, ymax=85, ylabel={failing actions executed unchecked (\% of episodes)}]
\addplot[pAV] table[x=calls_avoided_pct, y=unsafe_pct] {data/plane_av.dat};
\addplot[pNV] table[x=calls_avoided_pct, y=unsafe_pct] {data/plane_nv.dat};
\addplot[pRand] table[x=calls_avoided_pct, y=unsafe_pct] {data/plane_rand.dat};
\addplot[pObs] table[x=calls_avoided_pct, y=unsafe_pct] {data/plane_obs.dat};
\addplot[pComb] table[x=calls_avoided_pct, y=unsafe_pct] {data/plane_comb.dat};
\addplot[pInt] table[x=calls_avoided_pct, y=unsafe_pct] {data/plane_int.dat};
\addplot[cRzero, thick, -{Latex[length=2mm]}, quiver={u=\thisrow{x1}-\thisrow{x0}, v=\thisrow{y1}-\thisrow{y0}}]
  table[x=x0, y=y0] {data/plane_arrows.dat};
\node[font=\scriptsize, text=cRzero!80!black, align=center] at (axis cs:20,32) {arrows:\\R0 rule};
\draw[black!40, densely dashed] (axis cs:-2,-2) rectangle (axis cs:102,4);
\nextgroupplot[title={\footnotesize (b) zoom: 0--4\% unsafe}, ymin=-0.2, ymax=4,
  ylabel={}, ytick={0,1,2,3,4},
  legend style={at={(0.5,-0.3)}, anchor=north, legend columns=4, legend cell align=left, /tikz/every even column/.append style={column sep=6pt}},
  legend to name=planelegend]
\addplot[pAV] table[x=calls_avoided_pct, y=unsafe_pct] {data/plane_av.dat};\addlegendentry{always validate}
\addplot[pNV] table[x=calls_avoided_pct, y=unsafe_pct] {data/plane_nv.dat};\addlegendentry{never validate}
\addplot[pRand] table[x=calls_avoided_pct, y=unsafe_pct] {data/plane_rand.dat};\addlegendentry{random (matched rates)}
\addplot[pObs] table[x=calls_avoided_pct, y=unsafe_pct] {data/plane_obs.dat};\addlegendentry{observables}
\addplot[pComb] table[x=calls_avoided_pct, y=unsafe_pct] {data/plane_comb.dat};\addlegendentry{obs.\,+\,internals\,+\,grounding}
\addplot[pInt] table[x=calls_avoided_pct, y=unsafe_pct] {data/plane_int.dat};\addlegendentry{internals only}
\addplot[pR0] table[x=calls_avoided_pct, y=unsafe_pct] {data/plane_comb_r0.dat};\addlegendentry{R0 variant (ring)}
\addplot[pR0, forget plot] table[x=calls_avoided_pct, y=unsafe_pct] {data/plane_int_r0.dat};
\addplot[pR0, forget plot] table[x=calls_avoided_pct, y=unsafe_pct] {data/plane_obs_r0.dat};
\addplot[pComb, forget plot] table[x=calls_avoided_pct, y=unsafe_pct] {data/plane_comb_r0.dat};
\addplot[pInt, forget plot] table[x=calls_avoided_pct, y=unsafe_pct] {data/plane_int_r0.dat};
\addplot[pObs, forget plot] table[x=calls_avoided_pct, y=unsafe_pct] {data/plane_obs_r0.dat};
\end{groupplot}
\node[anchor=north] at ($(group c1r1.south)!0.5!(group c2r1.south)+(0,-0.95)$) {\pgfplotslegendfromname{planelegend}};
\end{tikzpicture}

\caption{Closed-loop safety--savings plane (400 fresh CSTR episodes per model; one marker per model and
policy). Ideal policies lie at the bottom right. Random routing at the screens' rates executes many failing
actions; trained screens skip as many calls with few failures. Two internals-based screens (Qwen-7B combined,
SmolLM2 combined and internals-only) accepted failing \emph{retries}; the arrows show the effect of the
first-proposal-only rule R0, which moves them to the bottom of the plane. Panel~(b) enlarges the dashed region of
panel~(a); rings mark R0 variants.}
\label{fig:plane}
\end{figure}

\begin{table*}[!tbp]
\centering
\small
\caption{Closed-loop CSTR recovery on 400 fresh episodes per model. $\Delta$: paired difference to always-validate (percentage points, bootstrap 95\% CI over episodes). Unsafe: failing actions executed without validation. Calls: validator calls avoided relative to always-validate. R0: the probe may skip validation only for first proposals (retries are always validated).}
\label{tab:closedloop}
\resizebox{\textwidth}{!}{%
\begin{tabular}{llrrrr}
\toprule
Model & Policy & Recovered (\%) & $\Delta$ recovered & Unsafe & Calls avoided \\
\midrule
Qwen-1.5B & Always validate & 34.5 & -- & 0 & -- \\
 & Observables & 33.0 & -1.5 [-3.0, -0.2] & 12 & 67\% \\
 & Obs.+int.+grd. & 33.2 & -1.2 [-2.5, +0.0] & 9 & 64\% \\
 & Internals only & 33.2 & -1.2 [-2.5, -0.2] & 0 & 35\% \\
 & Random (matched) & 33.8 & -0.8 [-1.8, +0.0] & 186 & 60\% \\
 & Never validate & 34.5 & +0.0 [+0.0, +0.0] & 259 & 100\% \\
\midrule
Qwen-3B & Always validate & 43.8 & -- & 0 & -- \\
 & Observables & 41.0 & -2.8 [-5.8, +0.2] & 0 & 47\% \\
 & Obs.+int.+grd. & 37.0 & -6.8 [-10.2, -3.5] & 1 & 72\% \\
 & Internals only & 30.0 & -13.8 [-17.5, -10.5] & 0 & 65\% \\
 & Random (matched) & 37.8 & -6.0 [-9.0, -3.0] & 110 & 47\% \\
 & Never validate & 30.0 & -13.8 [-17.2, -10.5] & 274 & 100\% \\
\midrule
Qwen-7B & Always validate & 34.0 & -- & 0 & -- \\
 & Observables & 25.0 & -9.0 [-12.8, -5.2] & 8 & 90\% \\
 & Observables, R0 & 25.5 & -8.5 [-12.2, -4.8] & 4 & 87\% \\
 & Obs.+int.+grd. & 26.5 & -7.5 [-11.0, -4.0] & 165 & 75\% \\
 & Obs.+int.+grd., R0 & 32.5 & -1.5 [-4.2, +1.2] & 4 & 41\% \\
 & Internals only & 30.0 & -4.0 [-6.2, -1.8] & 0 & 44\% \\
 & Internals only, R0 & 30.0 & -4.0 [-6.2, -1.8] & 0 & 44\% \\
 & Random (matched) & 19.2 & -14.8 [-18.5, -10.8] & 107 & 78\% \\
 & Never validate & 19.2 & -14.8 [-18.5, -11.2] & 318 & 100\% \\
\midrule
Llama-3B & Always validate & 53.8 & -- & 0 & -- \\
 & Observables & 35.2 & -18.5 [-24.2, -13.2] & 0 & 62\% \\
 & Obs.+int.+grd. & 44.5 & -9.2 [-14.2, -4.8] & 0 & 45\% \\
 & Internals only & 49.8 & -4.0 [-8.0, +0.0] & 0 & 29\% \\
 & Random (matched) & 29.5 & -24.2 [-29.5, -19.5] & 0 & 68\% \\
 & Never validate & 15.2 & -38.5 [-43.5, -33.8] & 312 & 100\% \\
\midrule
SmolLM-1.7B & Always validate & 36.5 & -- & 0 & -- \\
 & Observables & 35.2 & -1.2 [-2.2, -0.2] & 13 & 62\% \\
 & Observables, R0 & 35.2 & -1.2 [-2.2, -0.2] & 13 & 62\% \\
 & Obs.+int.+grd. & 35.0 & -1.5 [-2.8, -0.5] & 237 & 66\% \\
 & Obs.+int.+grd., R0 & 35.5 & -1.0 [-2.0, -0.2] & 11 & 65\% \\
 & Internals only & 35.5 & -1.0 [-2.0, -0.2] & 243 & 44\% \\
 & Internals only, R0 & 36.0 & -0.5 [-1.2, +0.0] & 2 & 42\% \\
 & Random (matched) & 34.0 & -2.5 [-4.2, -1.0] & 172 & 65\% \\
 & Never validate & 36.0 & -0.5 [-1.2, +0.0] & 254 & 100\% \\
\bottomrule
\end{tabular}}
\end{table*}


\subsection{A failure mode: internals-based probes on retries}
Two models showed a failure that offline evaluation could not reveal. The combined screen of Qwen2.5-7B executed
165 failing actions, 161 of them on retries. SmolLM2's combined and internals-only screens executed 237 and 243,
nearly all on retries. The observable screens of the same models were largely unaffected: they executed 8 and
13 failing actions, and only 4 of Qwen2.5-7B's 8 were on retries.
Remark~\ref{rem:retries} explains this. The probes were fitted and certified on first proposals, and a retry
prompt contains feedback that a first prompt never contains. The internal state that the probe reads therefore
shifts: the same forward-pass statistics mean something different after the model has been told it failed. The
probe misreads failing retries as safe. Observable features of the plant state and the proposed change do not
depend on the feedback in the same way, so the observable probes largely kept their calibration. For the other three models the probes
accepted no failing retries, so the shift was present but harmless there.

\paragraph{The first-proposal-only rule restores the guarantee's premise} Under R0 a retry that the probe would
accept is validated instead (Algorithm~\ref{alg:episode}). Every unchecked execution is then a first proposal,
which is the population covered by Proposition~\ref{prop:cert}. The rule was pre-registered after the failure
was observed and before any run with it, and it was run for the two affected models.
\begin{itemize}
\item \textbf{SmolLM2:} failing executions fell from 237 to 11 (combined) and from 243 to 2 (internals only).
Recovery was within 1 point of always-validate ($-1.0$ [$-2.0$, $-0.2$] and $-0.5$ [$-1.2$, $0.0$]), and 65\% and
42\% of validator calls were still avoided.
\item \textbf{Qwen2.5-7B:} failing executions of the combined screen fell from 165 to 4. Recovery was unchanged
within uncertainty ($-1.5$ [$-4.2$, $+1.2$] points) while 41\% of calls were avoided. Its internals-only screen had
accepted no failing retry and was unchanged by the rule (0 failing executions, 44\% of calls avoided).
\end{itemize}
For the other three models R0 does not change safety. In 400 episodes their probes accepted only five retries
(Qwen2.5-3B, combined screen), and all five passed; R0 would have validated them at the cost of five extra
calls. With R0, the \textbf{internals-only screens executed 0, 0, 0, 0 and 2 failing actions unchecked in the
2000 closed-loop episodes} of the five models, while avoiding 35\%, 65\%, 44\%, 29\% and 42\% of validator calls.
This screen uses no plant measurement and no model of the plant. Its safety comes from the agent's own forward
pass and the first-proposal-only rule.

\subsection{Where recovery is lost}
The screens' recovery cost ranged from $-0.5$ to $-18.5$ points against always-validate. It came almost entirely
from unchecked \emph{rejects}, not from unsafe accepts. Each rejection consumes one of the six proposals and
returns feedback without measurements. Some rejected proposals were good: up to 0.45 per episode for Qwen2.5-3B's
internals-only screen, whose recovery fell by 13.8 points. Llama-3.2-3B's proposals rarely passed and none of its
screens accepted anything, so every skipped call was a rejection. Its recovery fell by 4.0 points (internals
only) to 18.5 points (observables). For the two smallest models validation did not improve recovery at all:
never-validate recovered as often as always-validate (34.5\% versus 34.5\% for Qwen2.5-1.5B, 36.0\% versus 36.5\%
for SmolLM2), because these models mostly repeat their proposals. For them the validator, and the screens, serve
only to prevent the execution of failing actions.

\subsection{Compute and the break-even validator cost}
Because the plant is paused during decisions and both decoding and validation are deterministic, the logged
decisions can be re-costed exactly for any validator cost (Section~\ref{sec:costmodel}).
Figure~\ref{fig:cost} plots Eq.~\eqref{eq:cost} for each model, using the R0 variant wherever it was run, and
Table~\ref{tab:costs} gives the savings and break-even costs. On our hardware generation dominated the
per-episode compute (17--167\,s per episode under always-validate), and the extra passes needed for internal
signals added 2--14\,s per episode. At the measured validator cost (8--12\,s per call), the screens saved little
or cost time. Every unchecked rejection triggers a regeneration that costs about as much as the validation it
replaces. As the validator becomes more expensive, the savings grow. Excluding Llama-3.2-3B, the deployed screens
saved 11--47\% of total compute at 30\,s per validation and 21--62\% at 60\,s, with break-even validator costs
of 3--18\,s per call. As $\cval$ grows, the relative saving approaches the fraction of calls avoided
(Section~\ref{sec:costmodel}). Llama-3.2-3B is the exception. It is screened only by rejection and generates
slowly, so its break-even lies at 73--107\,s, and its timings were further inflated by GPU memory pressure in our
runs.

\begin{figure}[!htbp]
\centering
% Compute per closed-loop episode as a linear function of the validator's cost per call (exact re-costing of the
% logged decisions). Dots: break-even with always-validate; dashed: measured validator cost. Data: data/cost_m*.dat.
\begin{tikzpicture}
\pgfplotsset{
  cAV/.style={black, very thick},
  cO/.style={cObs, thick},
  cC/.style={cComb, thick, densely dashed},
  cI/.style={cInt, thick, dash dot},
  costpanel/.style={title={\footnotesize #1}},
}
\newcommand{\costlines}[1]{%
  \addplot[black!45, densely dashed, forget plot] table[x=c, y=C] {data/measured_m#1.dat};
  \addplot[cAV] table[x=c, y=av] {data/cost_m#1.dat};
  \addplot[cO] table[x=c, y=obs] {data/cost_m#1.dat};
  \addplot[cC] table[x=c, y=comb] {data/cost_m#1.dat};
  \addplot[cI] table[x=c, y=int] {data/cost_m#1.dat};
  \addplot[only marks, mark=*, mark size=1.6pt, black, forget plot] table[x=c, y=C] {data/breakeven_m#1.dat};}
\begin{groupplot}[group style={group size=3 by 2, horizontal sep=11mm, vertical sep=21mm},
  paperaxis, width=0.33\textwidth, height=0.3\textwidth, xmin=0, xmax=60, ymin=0, xtick={0,20,40,60},
  xlabel={$c_{\mathrm{val}}$ (s per call)}, xmajorgrids]
\nextgroupplot[costpanel={Qwen-1.5B}, ylabel={compute per episode (s)}]\costlines{1}
\nextgroupplot[costpanel={Qwen-3B}]\costlines{2}
\nextgroupplot[costpanel={Qwen-7B (4-bit)}]\costlines{3}
\nextgroupplot[costpanel={Llama-3B}, ylabel={compute per episode (s)}]\costlines{4}
\nextgroupplot[costpanel={SmolLM-1.7B}]\costlines{5}
\nextgroupplot[hide axis, xmin=0, xmax=1, ymin=0, ymax=1,
  legend style={at={(0.5,0.5)}, anchor=center, draw=black!30, fill=white, font=\scriptsize}, legend cell align=left]
\addlegendimage{cAV}\addlegendentry{always validate}
\addlegendimage{cO}\addlegendentry{observables}
\addlegendimage{cC}\addlegendentry{obs.\,+\,internals\,+\,grounding$^{\ast}$}
\addlegendimage{cI}\addlegendentry{internals only$^{\ast}$}
\addlegendimage{only marks, mark=*, mark size=1.6pt, black}\addlegendentry{break-even}
\addlegendimage{black!45, densely dashed}\addlegendentry{measured $c_{\mathrm{val}}$}
\end{groupplot}
\end{tikzpicture}

\caption{Compute per closed-loop episode as a function of the validator's cost per call, from exact re-costing of
the logged decisions (Eq.~\ref{eq:cost}). Dots mark the break-even cost with always-validate (Eq.~\ref{eq:breakeven});
the dashed line marks the measured validator cost. $^{\ast}$First-proposal-only (R0) variant for Qwen-7B and
SmolLM2. For the other three models R0 was not run separately, because it does not change their unchecked
executions. The screen's cost is fixed, whereas the validator's grows with simulation horizon and fidelity, so
the advantage of screening widens to the right.}
\label{fig:cost}
\end{figure}

\begin{table*}[!tbp]
\centering
\small
\caption{Compute per episode relative to always-validate (positive = saving). The logged closed-loop decisions do not depend on validator latency (the plant is paused during decisions), so the same decisions are re-costed exactly for a validator taking 30\,s or 60\,s per call. Break-even: validator cost per call above which the policy saves time.}
\label{tab:costs}
\resizebox{\textwidth}{!}{%
\begin{tabular}{llrrrr}
\toprule
Model & Policy & Measured & 30\,s & 60\,s & Break-even (s) \\
\midrule
Qwen-1.5B & Observables probe & +11\% & +28\% & +40\% & 3 \\
 & Obs. + internals + grounding probe & +5\% & +23\% & +36\% & 7 \\
 & Internals-only probe & +6\% & +14\% & +21\% & 3 \\
\midrule
Qwen-3B & Observables probe & +2\% & +21\% & +30\% & 7 \\
 & Obs. + internals + grounding probe & -4\% & +28\% & +44\% & 10 \\
 & Internals-only probe & -19\% & +16\% & +34\% & 17 \\
\midrule
Llama-3B & Observables probe & -40\% & -22\% & -5\% & 73 \\
 & Obs. + internals + grounding probe & -40\% & -24\% & -11\% & 97 \\
 & Internals-only probe & -29\% & -18\% & -9\% & 107 \\
\midrule
Qwen-7B & Observables probe & +18\% & +49\% & +64\% & 2 \\
 & Observables probe, first-proposal-only & +17\% & +47\% & +62\% & 3 \\
 & Obs. + internals + grounding probe & +31\% & +50\% & +59\% & any \\
 & Obs. + internals + grounding, first-proposal-only & -1\% & +17\% & +26\% & 9 \\
 & Internals-only probe & -1\% & +18\% & +28\% & 9 \\
 & Internals-only, first-proposal-only & -2\% & +17\% & +27\% & 10 \\
\midrule
SmolLM-1.7B & Observables probe & +0\% & +24\% & +39\% & 12 \\
 & Observables probe, first-proposal-only & -7\% & +20\% & +36\% & 15 \\
 & Obs. + internals + grounding probe & +44\% & +52\% & +58\% & any \\
 & Obs. + internals + grounding, first-proposal-only & -1\% & +25\% & +41\% & 12 \\
 & Internals-only probe & +40\% & +42\% & +43\% & any \\
 & Internals-only, first-proposal-only & -9\% & +11\% & +23\% & 18 \\
\bottomrule
\end{tabular}}
\end{table*}

